% --- Archon physics formalization source begin ---
% archon:physics
% archon:covers PhyXMiniProblems/problem_phyx_mini_0613.lean
% archon:source-report reports/phyx_mini/problem_phyx_mini_0613.source.json

\chapter{Physics problem phyx\_mini\_0613}
\label{ch:PhyXMiniProblems_problem_phyx_mini_0613}

\paragraph{Problem source.}
\#\# Physical scenario

A photon with wavelength \$\textbackslash{}lambda\$ is incident on a stationary particle with mass \$M\$, as shown in figure. The photon is annihilated while an electron–positron pair is produced. The target particle moves off in the original direction of the photon with speed \$V\_M\$. The electron travels with speed \$v\$ at angle \$\textbackslash{}phi\$ with respect to that direction. Owing to momentum conservation, the positron has the same speed as the electron.

\#\# Question

What is the energy of the incident photon in this case?

\#\# Answer choices

- A: A:  10.0MeV

- B: B:  11.2MeV

- C: C:  12.4MeV

- D: D:  13.3MeV

\#\# Figure caption (auxiliary; use the image as primary evidence)

The image is divided into two sections labeled "BEFORE" and "AFTER."

**Before:**
- On the left, there is a wavy purple arrow labeled with the Greek letter λ.
- The arrow is pointing towards an orange circle marked with the letter "M."

**After:**
- On the right, there is a setup featuring three different labeled objects.
- The orange circle labeled "M" is present again.
- A blue circle labeled "e⁻" (representing an electron) and a red circle labeled "e⁺" (representing a positron) are positioned below and above, respectively.
- Green arrows labeled "ν" are pointing outward from both the blue and red circles.
- A solid green arrow labeled "Vₘ" extends from the orange circle "M."
- Dashed lines, labeled with the Greek letter "φ," connect "M" to the paths between "e⁻" and "e⁺."

Overall, the diagram appears to illustrate a transformation or interaction involving the elements "M," "e⁻," and "e⁺," along with associated particles or forces.

\#\# Dataset metadata

Quantum Phenomena; Physical Model Grounding Reasoning, Multi-Formula Reasoning

\paragraph{Recorded answer/context.}
A: A:  10.0MeV

\paragraph{Figure/image path.}
/root/proposal\_for\_physic/science-mango/phyx\_mini\_run/phyx\_data/test\_image/613.png

\paragraph{Formalization target.}
create a compiling Lean file with sorry bodies at `PhyXMiniProblems/problem\_phyx\_mini\_0613.lean`. The Lean declarations must preserve the physical quantities, dimensions or dimensional roles, figure labels, governing-law hypotheses, and final relation expressed by this problem.
Use Mathlib/Physlib names found through LeanExplore where available. If a physics API is missing, introduce faithful local abstractions rather than scalar placeholder aliases.

\begin{theorem}[Physics formalization target]
\label{thm:physics:phyx_mini_0613:target}
\lean{PhyXMiniProblems.ProblemPhyXMini0613.problem_phyx_mini_0613}
\uses{def:physics:phyx-mini-0613:phyxminiproblems-problemphyxmini0613-wavelengthinmeters, def:physics:phyx-mini-0613:phyxminiproblems-problemphyxmini0613-energyinjoules, def:physics:phyx-mini-0613:phyxminiproblems-problemphyxmini0613-speedoflightinmeterspersecond, def:physics:phyx-mini-0613:phyxminiproblems-problemphyxmini0613-speedratiotolight, def:physics:phyx-mini-0613:phyxminiproblems-problemphyxmini0613-restenergyinjoules, def:physics:phyx-mini-0613:phyxminiproblems-problemphyxmini0613-photonpairproductionsetup, def:physics:phyx-mini-0613:phyxminiproblems-problemphyxmini0613-matchespairproductionscenario, def:physics:phyx-mini-0613:phyxminiproblems-problemphyxmini0613-matchesprimarypairproductionfigure, def:physics:phyx-mini-0613:phyxminiproblems-problemphyxmini0613-hasphysicalpairproductionparameters, def:physics:phyx-mini-0613:phyxminiproblems-problemphyxmini0613-satisfiesrelativisticpairproductionlaws}
The assigned autoformalize agent should translate this physics problem into Lean declarations in the covered file, with theorem and lemma proof bodies written as `by sorry`.
\end{theorem}
\begin{proof}
This is an autoformalization task, not a proof task. Produce statements that can later be proved by the physics prover without weakening the physical model.
\end{proof}

% --- Archon declaration topology begin ---
\section{Declaration topology}
\begin{definition}[Rest Mass Quantity]
  \label{def:physics:phyx-mini-0613:phyxminiproblems-problemphyxmini0613-restmassquantity}
  \lean{PhyXMiniProblems.ProblemPhyXMini0613.RestMassQuantity}
  A nonnegative, unit-independent physical rest mass
\end{definition}
\begin{proof}
  This is the declared model definition of Rest Mass Quantity.
\end{proof}
\begin{definition}[Wavelength Quantity]
  \label{def:physics:phyx-mini-0613:phyxminiproblems-problemphyxmini0613-wavelengthquantity}
  \lean{PhyXMiniProblems.ProblemPhyXMini0613.WavelengthQuantity}
  A nonnegative, unit-independent physical wavelength
\end{definition}
\begin{proof}
  This is the declared model definition of Wavelength Quantity.
\end{proof}
\begin{definition}[Plane Momentum]
  \label{def:physics:phyx-mini-0613:phyxminiproblems-problemphyxmini0613-planemomentum}
  \lean{PhyXMiniProblems.ProblemPhyXMini0613.PlaneMomentum}
  A dimensionful two-dimensional momentum in the plane of the figure
\end{definition}
\begin{proof}
  This is the declared model definition of Plane Momentum.
\end{proof}
\begin{definition}[Diagram Axis]
  \label{def:physics:phyx-mini-0613:phyxminiproblems-problemphyxmini0613-diagramaxis}
  \lean{PhyXMiniProblems.ProblemPhyXMini0613.DiagramAxis}
  The two coordinate axes determined by the incident-photon direction
\end{definition}
\begin{proof}
  This is the declared model definition of Diagram Axis.
\end{proof}
\begin{definition}[to Fin]
  \label{def:physics:phyx-mini-0613:phyxminiproblems-problemphyxmini0613-diagramaxis-tofin}
  \lean{PhyXMiniProblems.ProblemPhyXMini0613.DiagramAxis.toFin}
  \uses{def:physics:phyx-mini-0613:phyxminiproblems-problemphyxmini0613-diagramaxis}
  Coordinate index corresponding to a named diagram axis
\end{definition}
\begin{proof}
  This is the declared model definition of to Fin.
\end{proof}
\begin{definition}[rest Mass In Kilograms]
  \label{def:physics:phyx-mini-0613:phyxminiproblems-problemphyxmini0613-restmassinkilograms}
  \lean{PhyXMiniProblems.ProblemPhyXMini0613.restMassInKilograms}
  \uses{def:physics:phyx-mini-0613:phyxminiproblems-problemphyxmini0613-restmassquantity}
  SI mass readout in kilograms
\end{definition}
\begin{proof}
  This is the declared model definition of rest Mass In Kilograms.
\end{proof}
\begin{definition}[wavelength In Meters]
  \label{def:physics:phyx-mini-0613:phyxminiproblems-problemphyxmini0613-wavelengthinmeters}
  \lean{PhyXMiniProblems.ProblemPhyXMini0613.wavelengthInMeters}
  \uses{def:physics:phyx-mini-0613:phyxminiproblems-problemphyxmini0613-wavelengthquantity}
  SI wavelength readout in metres
\end{definition}
\begin{proof}
  This is the declared model definition of wavelength In Meters.
\end{proof}
\begin{definition}[energy In Joules]
  \label{def:physics:phyx-mini-0613:phyxminiproblems-problemphyxmini0613-energyinjoules}
  \lean{PhyXMiniProblems.ProblemPhyXMini0613.energyInJoules}
  SI energy readout in joules
\end{definition}
\begin{proof}
  This is the declared model definition of energy In Joules.
\end{proof}
\begin{definition}[energy In Mega Electron Volts]
  \label{def:physics:phyx-mini-0613:phyxminiproblems-problemphyxmini0613-energyinmegaelectronvolts}
  \lean{PhyXMiniProblems.ProblemPhyXMini0613.energyInMegaElectronVolts}
  \uses{def:physics:phyx-mini-0613:phyxminiproblems-problemphyxmini0613-energyinjoules}
  Energy readout in megaelectron-volts
\end{definition}
\begin{proof}
  This is the declared model definition of energy In Mega Electron Volts.
\end{proof}
\begin{definition}[speed In Meters Per Second]
  \label{def:physics:phyx-mini-0613:phyxminiproblems-problemphyxmini0613-speedinmeterspersecond}
  \lean{PhyXMiniProblems.ProblemPhyXMini0613.speedInMetersPerSecond}
  SI speed readout in metres per second
\end{definition}
\begin{proof}
  This is the declared model definition of speed In Meters Per Second.
\end{proof}
\begin{definition}[speed Of Light In Meters Per Second]
  \label{def:physics:phyx-mini-0613:phyxminiproblems-problemphyxmini0613-speedoflightinmeterspersecond}
  \lean{PhyXMiniProblems.ProblemPhyXMini0613.speedOfLightInMetersPerSecond}
  Physlib's exact vacuum speed of light, read in metres per second
\end{definition}
\begin{proof}
  This is the declared model definition of speed Of Light In Meters Per Second.
\end{proof}
\begin{definition}[speed Ratio To Light]
  \label{def:physics:phyx-mini-0613:phyxminiproblems-problemphyxmini0613-speedratiotolight}
  \lean{PhyXMiniProblems.ProblemPhyXMini0613.speedRatioToLight}
  \uses{def:physics:phyx-mini-0613:phyxminiproblems-problemphyxmini0613-speedinmeterspersecond, def:physics:phyx-mini-0613:phyxminiproblems-problemphyxmini0613-speedoflightinmeterspersecond}
  A massive particle's speed as the dimensionless ratio v/c
\end{definition}
\begin{proof}
  This is the declared model definition of speed Ratio To Light.
\end{proof}
\begin{definition}[momentum Component In SI]
  \label{def:physics:phyx-mini-0613:phyxminiproblems-problemphyxmini0613-momentumcomponentinsi}
  \lean{PhyXMiniProblems.ProblemPhyXMini0613.momentumComponentInSI}
  \uses{def:physics:phyx-mini-0613:phyxminiproblems-problemphyxmini0613-planemomentum, def:physics:phyx-mini-0613:phyxminiproblems-problemphyxmini0613-diagramaxis}
  One plane-momentum component in coherent SI units kg m / s
\end{definition}
\begin{proof}
  This is the declared model definition of momentum Component In SI.
\end{proof}
\begin{definition}[momentum Magnitude In SI]
  \label{def:physics:phyx-mini-0613:phyxminiproblems-problemphyxmini0613-momentummagnitudeinsi}
  \lean{PhyXMiniProblems.ProblemPhyXMini0613.momentumMagnitudeInSI}
  \uses{def:physics:phyx-mini-0613:phyxminiproblems-problemphyxmini0613-planemomentum, def:physics:phyx-mini-0613:phyxminiproblems-problemphyxmini0613-momentumcomponentinsi}
  Euclidean magnitude of a plane momentum in coherent SI units
\end{definition}
\begin{proof}
  This is the declared model definition of momentum Magnitude In SI.
\end{proof}
\begin{definition}[rest Energy In Joules]
  \label{def:physics:phyx-mini-0613:phyxminiproblems-problemphyxmini0613-restenergyinjoules}
  \lean{PhyXMiniProblems.ProblemPhyXMini0613.restEnergyInJoules}
  \uses{def:physics:phyx-mini-0613:phyxminiproblems-problemphyxmini0613-restmassquantity, def:physics:phyx-mini-0613:phyxminiproblems-problemphyxmini0613-restmassinkilograms, def:physics:phyx-mini-0613:phyxminiproblems-problemphyxmini0613-speedoflightinmeterspersecond}
  Rest-energy readout m c\textasciicircum{}2 in joules
\end{definition}
\begin{proof}
  This is the declared model definition of rest Energy In Joules.
\end{proof}
\begin{definition}[Massive State Label]
  \label{def:physics:phyx-mini-0613:phyxminiproblems-problemphyxmini0613-massivestatelabel}
  \lean{PhyXMiniProblems.ProblemPhyXMini0613.MassiveStateLabel}
  The four massive-particle states before and after the interaction
\end{definition}
\begin{proof}
  This is the declared model definition of Massive State Label.
\end{proof}
\begin{definition}[Reaction Channel]
  \label{def:physics:phyx-mini-0613:phyxminiproblems-problemphyxmini0613-reactionchannel}
  \lean{PhyXMiniProblems.ProblemPhyXMini0613.ReactionChannel}
  The reaction channel described in the problem, distinguished from others
\end{definition}
\begin{proof}
  This is the declared model definition of Reaction Channel.
\end{proof}
\begin{definition}[Massive Particle State]
  \label{def:physics:phyx-mini-0613:phyxminiproblems-problemphyxmini0613-massiveparticlestate}
  \lean{PhyXMiniProblems.ProblemPhyXMini0613.MassiveParticleState}
  \uses{def:physics:phyx-mini-0613:phyxminiproblems-problemphyxmini0613-restmassquantity, def:physics:phyx-mini-0613:phyxminiproblems-problemphyxmini0613-planemomentum}
  A massive particle's independent relativistic state variables
\end{definition}
\begin{proof}
  This is the declared model definition of Massive Particle State.
\end{proof}
\begin{definition}[Photon State]
  \label{def:physics:phyx-mini-0613:phyxminiproblems-problemphyxmini0613-photonstate}
  \lean{PhyXMiniProblems.ProblemPhyXMini0613.PhotonState}
  \uses{def:physics:phyx-mini-0613:phyxminiproblems-problemphyxmini0613-wavelengthquantity, def:physics:phyx-mini-0613:phyxminiproblems-problemphyxmini0613-planemomentum}
  The incident photon's independent physical state variables
\end{definition}
\begin{proof}
  This is the declared model definition of Photon State.
\end{proof}
\begin{definition}[Figure Panel]
  \label{def:physics:phyx-mini-0613:phyxminiproblems-problemphyxmini0613-figurepanel}
  \lean{PhyXMiniProblems.ProblemPhyXMini0613.FigurePanel}
  The two panels printed in the primary image
\end{definition}
\begin{proof}
  This is the declared model definition of Figure Panel.
\end{proof}
\begin{definition}[Figure Entity]
  \label{def:physics:phyx-mini-0613:phyxminiproblems-problemphyxmini0613-figureentity}
  \lean{PhyXMiniProblems.ProblemPhyXMini0613.FigureEntity}
  Object labels visible in one or both panels
\end{definition}
\begin{proof}
  This is the declared model definition of Figure Entity.
\end{proof}
\begin{definition}[Figure Quantity Label]
  \label{def:physics:phyx-mini-0613:phyxminiproblems-problemphyxmini0613-figurequantitylabel}
  \lean{PhyXMiniProblems.ProblemPhyXMini0613.FigureQuantityLabel}
  Symbolic quantity labels printed beside arrows or angle marks
\end{definition}
\begin{proof}
  This is the declared model definition of Figure Quantity Label.
\end{proof}
\begin{definition}[Planar Direction]
  \label{def:physics:phyx-mini-0613:phyxminiproblems-problemphyxmini0613-planardirection}
  \lean{PhyXMiniProblems.ProblemPhyXMini0613.PlanarDirection}
  Qualitative arrow directions in the primary image
\end{definition}
\begin{proof}
  This is the declared model definition of Planar Direction.
\end{proof}
\begin{definition}[Pair Production Figure]
  \label{def:physics:phyx-mini-0613:phyxminiproblems-problemphyxmini0613-pairproductionfigure}
  \lean{PhyXMiniProblems.ProblemPhyXMini0613.PairProductionFigure}
  \uses{def:physics:phyx-mini-0613:phyxminiproblems-problemphyxmini0613-figurepanel, def:physics:phyx-mini-0613:phyxminiproblems-problemphyxmini0613-figureentity, def:physics:phyx-mini-0613:phyxminiproblems-problemphyxmini0613-figurequantitylabel, def:physics:phyx-mini-0613:phyxminiproblems-problemphyxmini0613-planardirection}
  Literal panel, label, and arrow information carried by the primary image
\end{definition}
\begin{proof}
  This is the declared model definition of Pair Production Figure.
\end{proof}
\begin{definition}[Photon Pair Production Setup]
  \label{def:physics:phyx-mini-0613:phyxminiproblems-problemphyxmini0613-photonpairproductionsetup}
  \lean{PhyXMiniProblems.ProblemPhyXMini0613.PhotonPairProductionSetup}
  \uses{def:physics:phyx-mini-0613:phyxminiproblems-problemphyxmini0613-massivestatelabel, def:physics:phyx-mini-0613:phyxminiproblems-problemphyxmini0613-reactionchannel, def:physics:phyx-mini-0613:phyxminiproblems-problemphyxmini0613-massiveparticlestate, def:physics:phyx-mini-0613:phyxminiproblems-problemphyxmini0613-photonstate, def:physics:phyx-mini-0613:phyxminiproblems-problemphyxmini0613-pairproductionfigure}
  The complete pair-production setup, with the answer energy left independent
\end{definition}
\begin{proof}
  This is the declared model definition of Photon Pair Production Setup.
\end{proof}
\begin{definition}[Matches Pair Production Scenario]
  \label{def:physics:phyx-mini-0613:phyxminiproblems-problemphyxmini0613-matchespairproductionscenario}
  \lean{PhyXMiniProblems.ProblemPhyXMini0613.MatchesPairProductionScenario}
  \uses{def:physics:phyx-mini-0613:phyxminiproblems-problemphyxmini0613-diagramaxis, def:physics:phyx-mini-0613:phyxminiproblems-problemphyxmini0613-speedinmeterspersecond, def:physics:phyx-mini-0613:phyxminiproblems-problemphyxmini0613-momentumcomponentinsi, def:physics:phyx-mini-0613:phyxminiproblems-problemphyxmini0613-photonpairproductionsetup}
  The prose scenario. The equality of the lepton speeds is explicitly supplied in the question. No field here fixes the incident photon energy
\end{definition}
\begin{proof}
  This is the declared model definition of Matches Pair Production Scenario.
\end{proof}
\begin{definition}[Matches Primary Pair Production Figure]
  \label{def:physics:phyx-mini-0613:phyxminiproblems-problemphyxmini0613-matchesprimarypairproductionfigure}
  \lean{PhyXMiniProblems.ProblemPhyXMini0613.MatchesPrimaryPairProductionFigure}
  \uses{def:physics:phyx-mini-0613:phyxminiproblems-problemphyxmini0613-momentumcomponentinsi, def:physics:phyx-mini-0613:phyxminiproblems-problemphyxmini0613-momentummagnitudeinsi, def:physics:phyx-mini-0613:phyxminiproblems-problemphyxmini0613-photonpairproductionsetup}
  Primary-image evidence. The component equations attach the printed phi to the upper electron and lower positron trajectories and attach the photon and target recoil arrows to the original positive horizontal direction
\end{definition}
\begin{proof}
  This is the declared model definition of Matches Primary Pair Production Figure.
\end{proof}
\begin{definition}[Has Physical Pair Production Parameters]
  \label{def:physics:phyx-mini-0613:phyxminiproblems-problemphyxmini0613-hasphysicalpairproductionparameters}
  \lean{PhyXMiniProblems.ProblemPhyXMini0613.HasPhysicalPairProductionParameters}
  \uses{def:physics:phyx-mini-0613:phyxminiproblems-problemphyxmini0613-restmassinkilograms, def:physics:phyx-mini-0613:phyxminiproblems-problemphyxmini0613-wavelengthinmeters, def:physics:phyx-mini-0613:phyxminiproblems-problemphyxmini0613-energyinjoules, def:physics:phyx-mini-0613:phyxminiproblems-problemphyxmini0613-speedinmeterspersecond, def:physics:phyx-mini-0613:phyxminiproblems-problemphyxmini0613-speedoflightinmeterspersecond, def:physics:phyx-mini-0613:phyxminiproblems-problemphyxmini0613-photonpairproductionsetup}
  Positivity and subluminality conditions selecting physical states
\end{definition}
\begin{proof}
  This is the declared model definition of Has Physical Pair Production Parameters.
\end{proof}
\begin{definition}[Satisfies Relativistic Pair Production Laws]
  \label{def:physics:phyx-mini-0613:phyxminiproblems-problemphyxmini0613-satisfiesrelativisticpairproductionlaws}
  \lean{PhyXMiniProblems.ProblemPhyXMini0613.SatisfiesRelativisticPairProductionLaws}
  \uses{def:physics:phyx-mini-0613:phyxminiproblems-problemphyxmini0613-diagramaxis, def:physics:phyx-mini-0613:phyxminiproblems-problemphyxmini0613-wavelengthinmeters, def:physics:phyx-mini-0613:phyxminiproblems-problemphyxmini0613-energyinjoules, def:physics:phyx-mini-0613:phyxminiproblems-problemphyxmini0613-speedinmeterspersecond, def:physics:phyx-mini-0613:phyxminiproblems-problemphyxmini0613-speedoflightinmeterspersecond, def:physics:phyx-mini-0613:phyxminiproblems-problemphyxmini0613-speedratiotolight, def:physics:phyx-mini-0613:phyxminiproblems-problemphyxmini0613-momentumcomponentinsi, def:physics:phyx-mini-0613:phyxminiproblems-problemphyxmini0613-momentummagnitudeinsi, def:physics:phyx-mini-0613:phyxminiproblems-problemphyxmini0613-restenergyinjoules, def:physics:phyx-mini-0613:phyxminiproblems-problemphyxmini0613-massivestatelabel, def:physics:phyx-mini-0613:phyxminiproblems-problemphyxmini0613-photonpairproductionsetup}
  Governing laws for the isolated relativistic interaction: |p\_gamma| = h / lambda, with h = 2 pi hbar; the photon obeys E\_gamma = c |p\_gamma|; every massive state obeys E = gamma(v/c) m c\textasciicircum{}2, E\textasciicircum{}2 = (m c\textasciicircum{}2)\textasciicircum{}2 + (p c)\textasciicircum{}2, and p c\textasciicircum{}2 = E v; total energy and both displayed momentum components are conserved. The photon laws are deliberately factored through momentum, so the requested symbolic energy--wavelength formula is a consequence rather than a premise.
\end{definition}
\begin{proof}
  This is the declared model definition of Satisfies Relativistic Pair Production Laws.
\end{proof}
\begin{lemma}[incident Photon Energy from energy Conservation]
  \label{lem:physics:phyx-mini-0613:phyxminiproblems-problemphyxmini0613-incidentphotonenergy-from-energyconservation}
  \lean{PhyXMiniProblems.ProblemPhyXMini0613.incidentPhotonEnergy_from_energyConservation}
  \uses{def:physics:phyx-mini-0613:phyxminiproblems-problemphyxmini0613-energyinjoules, def:physics:phyx-mini-0613:phyxminiproblems-problemphyxmini0613-photonpairproductionsetup, def:physics:phyx-mini-0613:phyxminiproblems-problemphyxmini0613-satisfiesrelativisticpairproductionlaws}
  The direct symbolic answer supplied by energy conservation, before inserting any missing numerical data
\end{lemma}
\begin{proof}
  The conclusion follows from the governing relations and figure readouts named in the statement of incident Photon Energy from energy Conservation.
\end{proof}
\begin{definition}[Answer Choice]
  \label{def:physics:phyx-mini-0613:phyxminiproblems-problemphyxmini0613-answerchoice}
  \lean{PhyXMiniProblems.ProblemPhyXMini0613.AnswerChoice}
  Labels printed beside the four answer choices
\end{definition}
\begin{proof}
  This is the declared model definition of Answer Choice.
\end{proof}
\begin{definition}[displayed Incident Photon Energy Me V]
  \label{def:physics:phyx-mini-0613:phyxminiproblems-problemphyxmini0613-displayedincidentphotonenergymev}
  \lean{PhyXMiniProblems.ProblemPhyXMini0613.displayedIncidentPhotonEnergyMeV}
  \uses{def:physics:phyx-mini-0613:phyxminiproblems-problemphyxmini0613-answerchoice}
  Incident-photon energy in MeV printed beside each answer choice
\end{definition}
\begin{proof}
  This is the declared model definition of displayed Incident Photon Energy Me V.
\end{proof}
\begin{definition}[recorded Answer Choice]
  \label{def:physics:phyx-mini-0613:phyxminiproblems-problemphyxmini0613-recordedanswerchoice}
  \lean{PhyXMiniProblems.ProblemPhyXMini0613.recordedAnswerChoice}
  \uses{def:physics:phyx-mini-0613:phyxminiproblems-problemphyxmini0613-answerchoice}
  Dataset metadata records answer choice A
\end{definition}
\begin{proof}
  This is the declared model definition of recorded Answer Choice.
\end{proof}
% --- Archon declaration topology end ---
% --- Archon physics formalization source end ---
